% !TEX root = ../main.tex
\note{Background: microsecond-scale RPC tail latency, per-core queues,
d-FCFS vs single-queue, and why queue length alone is a poor migration
signal.} Prior microsecond-scale schedulers~\cite{placeholder-shinjuku}
and work-conserving stealing~\cite{placeholder-workstealing} target work
conservation rather than per-request SLO outcomes.

RescueSched is not a generic load balancer. It does not aim to equalize
queue lengths. Instead, it treats cross-core migration as a scarce and
costly SLO-rescue action: a migration is useful only if it changes the
SLO outcome of a request without creating new target-side violations.
The central problem is:

\begin{quote}
Given a commodity multi-core RPC server with per-core request queues and
no preemption of running requests, can we reduce the SLO violation ratio
by selectively migrating only queued-but-not-running requests that are
locally doomed, remotely rescuable, and safe for the target core?
\end{quote}

\textbf{Scope and non-goals.} We target a single latency-critical RPC
service on one multi-core server with $N$ pinned worker cores, one FIFO
queue per core, microsecond-scale service times, request-level SLOs, no
preemption, no kernel modification, and no SmartNIC/programmable-switch
dependence. The scheduler migrates only queued request descriptors; it
never migrates running requests, application state, or payloads. We do
not attempt to solve global overload: if all cores are saturated with no
slack, RescueSched suppresses migration.